%杨舒云的实验报告编辑界面，使用了Huanyu Shi，2019级的模板，杨舒云在此拜谢ORZ!

%!TEX program = xelatex
\documentclass[dvipsnames, svgnames,a4paper,11pt]{article}
\input{YsySettings} 
\usepackage{lipsum}
\usepackage{adjustbox}

\begin{document}
	
	
	
	%实验报告封面	
	\begin{table}
		\renewcommand\arraystretch{1.7}
		\begin{tabularx}{\textwidth}{
				|X|X|X|X
				|X|X|X|X|}
			\hline
			\multicolumn{2}{|c|}{预习报告}&\multicolumn{2}{|c|}{实验记录}&\multicolumn{2}{|c|}{分析讨论}&\multicolumn{2}{|c|}{总成绩}\\
			\hline
			\LARGE25 & & \LARGE30 & & \LARGE25 & & \LARGE80 & \\
			\hline
		\end{tabularx}
	\end{table}
	
	\begin{table}
		\renewcommand\arraystretch{1.7}
		\begin{tabularx}{\textwidth}{|X|X|X|X|}
			\hline
			年级、专业： & 2022级 物理学 &组号： & 2\\
			\hline
			姓名： & 杨舒云  & 学号： & 22344020\\
			\hline
			实验时间： & 2023/11/30 & 教师签名： & \\
			\hline
		\end{tabularx}
	\end{table}
	
	\begin{center}
		\LARGE LabX4 \quad 伯努利原理实验
	\end{center}
	
	\textbf{【实验报告注意事项】}
	\begin{enumerate}
		\item 实验报告由三部分组成：
		\begin{enumerate}
			\item 预习报告：课前认真研读实验讲义，弄清实验原理；实验所需的仪器设备、用具及其使用、完成课前预习思考题；了解实验需要测量的物理量，并根据要求提前准备实验记录表格（可以参考实验报告模板，可以打印）。\textcolor{red}{\textbf{（25分）}}
			\item 实验记录：认真、客观记录实验条件、实验过程中的现象以及数据。实验记录请用珠笔或者钢笔书写并签名（\textcolor{red}{\textbf{用铅笔记录的被认为无效}}）。\textcolor{red}{\textbf{保持原始记录，包括写错删除部分，如因误记需要修改记录，必须按规范修改。}}（不得输入电脑打印，但可扫描手记后打印扫描件）；离开前请实验教师检查记录并签名。\textcolor{red}{\textbf{（30分）}}
			\item 数据处理及分析讨论：处理实验原始数据（学习仪器使用类型的实验除外），对数据的可靠性和合理性进行分析；按规范呈现数据和结果（图、表），包括数据、图表按顺序编号及其引用；分析物理现象（含回答实验思考题，写出问题思考过程，必要时按规范引用数据）；最后得出结论。\textcolor{red}{\textbf{（25分）}}
		\end{enumerate}
		\textbf{实验报告就是将预习报告、实验记录、和数据处理与分析合起来，加上本页封面。\textcolor{red}{（80分）}}
		\item 每次完成实验后的一周内交\textbf{实验报告}（特殊情况不能超过两周）。
	\end{enumerate}
	
%	\textbf{【实验安全注意事项】}	
%	\begin{enumerate}
%		\item 实验中\textcolor{red}{\textbf{禁止用手触摸极板，注意高压危险！}}
%		\item 喷油器每次喷1~2次即可，\textcolor{red}{\textbf{竖直放置、防止摔碎！}}
%		\item 计时时\textcolor{red}{\textbf{眼睛要平行刻度线！}} 
%	\end{enumerate}
	
	\textbf{【特别鸣谢及模板说明】}	
	
	感谢2019级学长石寰宇为本实验报告提供\LaTeX 模板。\textcolor{red}{\textbf{由于原实验报告模板缺少实验编号，为方便在电脑上整理，故添加自命名编号LabX4}}
	
	
	
	\clearpage
	\tableofcontents
	\clearpage
	
	
	
	%预习报告	
	\setcounter{section}{0}
	\section{LabX4 伯努利原理实验 \quad\heiti 预习报告}
	
	\subsection{实验目的}
	\begin{enumerate}
		\item 验证伯努利能量方程，加深对机械能守恒在流体中的应用理解；通过对水力现象的实验分析，理解能量转换特性；
		\item 学会各种水头的测量和计量方法，掌握流速、流量、压强等水力要素的实验量测技能；
		\item 观察并理解毕托管测速方法的原理。
	\end{enumerate}
	
	\subsection{实验内容}
	\begin{enumerate}
		\item 利用实验中设置的回路和测压水柱，测量不同流速条件下，各测压管的水头变化。
		\item 对测量结果进行速度分析，了解能量损耗，分析文丘里管的特性。
	\end{enumerate}
	
	\subsection{仪器用具}
	\begin{table}[htbp]
		\centering
		\renewcommand\arraystretch{1.6}
		% \setlength{\tabcolsep}{10mm}
		\begin{tabular}{p{0.05\textwidth}|p{0.20\textwidth}|p{0.05\textwidth}|p{0.5\textwidth}}
			\hline
			编号& 仪器用具名称 & 数量 &  主要参数（型号，测量范围，测量精度等） \\
			\hline
			1& 实验回路 & 1 & 包含铝底板、实验管、文丘里管、离心泵、流量计 \\
			
			2& 测压水柱 & 1 & 包含带刻度的测压管$\times13$，细软管（带止水夹）$\times13$ \\
			
			3& 磁铁 & 1 & 用于流量计读数切换 \\
			
			4& 调节阀 & 1 & -- \\
			
			5& 粗软管 & 1 & 备用 \\
			
			6& 气吹 & 1 & -- \\
			
			7& 水平泡 & 1 & -- \\
			
			8& DC可调电源 & 1 & -- \\
			\hline
		\end{tabular}
	\end{table}
	
	
	
	\subsection{原理概述}
	\begin{itemize}
		\item 不可压缩理想液体是指流体的密度在流动过程中不发生变化，且流体内部没有摩擦力的液体。考虑能量损耗的实际液体是指流体的密度在流动过程中可能发生变化，且流体内部存在摩擦力的液体。对于不可压缩理想液体，其沿流线的流动满足伯努利方程：
		
		\[p + \frac{1}{2} \rho v^2 + \rho gh = \text{常数}\]
		
		其中，\(p\) 是流体的压强，\(\rho\) 是流体的密度，\(v\) 是流体的速度，\(g\) 是重力加速度，\(h\) 是流体的高度。这个方程表明，流体在流动过程中，压强能、动能和重力势能之和保持不变。
		
		对于考虑能量损耗的实际液体，其沿流线的流动满足伯努利方程：
		
		\[p_1 + \frac{1}{2} \rho v_1^2 + \rho gh_1 = p_2 + \frac{1}{2} \rho v_2^2 + \rho gh_2 + \Delta E\]
		
		其中，\(\Delta E\) 是流体在流动过程中损失的能量，包括摩擦损失和局部损失。这个方程表明，流体在流动过程中，压强能、动能和重力势能之和不再保持不变，而是会减少一定的值。
		
		\begin{figure}[htbp]
			\centering
			\includegraphics[width=0.5\textwidth]{LabX4GraA1.png}
			\caption{伯努利原理示意图（参考文献[1]）}
			\label{fig:figA1}
		\end{figure}
		
		\item 流体的沿程阻力损失是指流体在流动过程中，由于流体内部的摩擦力和流体与管壁的摩擦力，导致流体的动能和压强能减少的现象。流体的沿程阻力损失与以下几个量有关：
		
		\begin{itemize}
			\item 流体的性质，如密度、粘度、可压缩性等；
			\item 流体的流速，流速越大，沿程阻力损失越大；
			\item 管道的长度，长度越长，沿程阻力损失越大；
			\item 管道的直径，直径越小，沿程阻力损失越大；
			\item 管道的粗糙度，粗糙度越大，沿程阻力损失越大。
		\end{itemize}
		
		流体的沿程阻力损失可以用以下公式定量计算：
		
		\[\Delta p = \lambda \frac{l}{d} \frac{1}{2} \rho v^2\]
		
		其中，\(\Delta p\) 是流体的沿程阻力损失，\(\lambda\) 是摩擦阻力系数，\(l\) 是管道的长度，\(d\) 是管道的直径，\(\rho\) 是流体的密度，\(v\) 是流体的流速。
		
		\item 层流是指流体的流动呈现出规则的层状结构，各层之间没有相互穿插的现象。湍流是指流体的流动呈现出无规则的涡旋结构，各层之间有相互穿插的现象。层流和湍流的区别主要有以下几点：
		
		\begin{itemize}
			\item 层流的流速分布是对称的，湍流的流速分布是不对称的；
			\item 层流的流速变化是连续的，湍流的流速变化是突变的；
			\item 层流的流动是稳定的，湍流的流动是不稳定的；
			\item 层流的阻力损失是小的，湍流的阻力损失是大的；
			\item 层流的流动是静音的，湍流的流动是嘈杂的。
		\end{itemize}
		
		\item 毕托管是一种用于测量流体流速的仪器，其工作原理是利用伯努利方程和连续方程。毕托管由一个细长的管子构成，管子的一端是封闭的，另一端是开口的，开口的部分与流体的流动方向垂直，称为\textbf{测压口}。测压口的内部有一个小孔，与管子的侧壁相连，称为\textbf{静压孔}。管子的侧壁还有另一个小孔，与流体的流动方向平行，称为\textbf{总压孔}。静压孔和总压孔分别与两个压力表相连，用于测量流体的静压和总压。
		
		\begin{figure}[htbp]
			\centering
			\includegraphics[width=0.5\textwidth]{LabX4GraA2.jpg}
			\caption{文丘里管示意图（参考文献[1]）}
			\label{fig:figA2}
		\end{figure}
		
		当流体流过毕托管时，由于测压口的阻挡，流体的速度会降为零，压强会升到最大，这个压强就是总压。由于静压孔的存在，流体的压强会从总压降到与周围流体相同的静压。由于总压孔的存在，流体的压强会保持为总压。因此，通过测量两个压力表的读数，就可以得到流体的总压和静压的差值，即\textbf{动压}。
		
		根据伯努利方程，流体的动压与流体的速度有如下关系：
		
		\[\frac{1}{2} \rho v^2 = p_0 - p\]
		
		其中，\(\rho\) 是流体的密度，\(v\) 是流体的速度，\(p_0\) 是流体的总压，\(p\) 是流体的静压。因此，通过测量动压，就可以计算出流体的速度。
		
		根据连续方程，流体的速度与流体的流量有如下关系：
		
		\[Q = Av\]
		
		其中，\(Q\) 是流体的流量，\(A\) 是流体的横截面积，\(v\) 是流体的速度。因此，通过测量速度，就可以计算出流体的流量。
	\end{itemize}
	
	
	\subsection{实验方案、步骤设计}
	\textbf{实验方案：}
	
	\begin{enumerate}
		\item 控制参数：泵的转速。通过调节泵的转速，可以改变实验管道中的液体流量，从而影响液体的流速和压强。控制范围：根据流量计的量程，选择合适的泵转速，使流量在0到10 L/min之间变化。
		\item 测量记录的物理量：流量、压强、流速。流量通过流量计直接读出，压强通过测压管的水位差计算，流速通过伯努利方程和连续方程计算。
		\item 设计实验记录表：实验记录表的目的是记录实验过程中的各种物理量，以便于后续的数据处理和分析。实验记录表应该包括以下内容：
		
		\begin{itemize}
			\item 实验条件：记录实验时的泵转速和流量，以及实验管道的直径和长度等参数；
			\item 实验数据：记录各测压管的水位，以及计算出的压强和流速；
			\item 实验结果：记录各种水头和总水头，以及计算出的能量损耗和摩擦压降等指标。
		\end{itemize}
		
		\begin{table}[h]
			\centering
			\caption{实验表格设计}
			\begin{tabular}{|c|c|c|c|c|c|c|c|c|c|}
				\hline
				流量设置值 $m^3/h$ & 0.0 & 0.1 & 0.15 & 0.2 & 0.25 & 0.3 & 0.35 & 0.4 & 0.45 \\
				\hline
				流量计流量 $m^3/h$ & & & & & & & & & \\
				\hline
				测压管水柱高度/mm & & & & & & & & & \\
				\hline
				1($\phi$20mm) & & & & & & & & & \\
				\hline
				2($\phi$20mm) & & & & & & & & & \\
				\hline
				3($\phi$15mm) & & & & & & & & & \\
				\hline
				4($\phi$15mm) & & & & & & & & & \\
				\hline
				5($\phi$15mm) & & & & & & & & & \\
				\hline
				6($\phi$15mm) & & & & & & & & & \\
				\hline
				7($\phi$15mm) & & & & & & & & & \\
				\hline
				8($\phi$15mm) & & & & & & & & & \\
				\hline
				9($\phi$15mm) & & & & & & & & & \\
				\hline
				10($\phi$15mm) & & & & & & & & & \\
				\hline
				11($\phi$15mm) & & & & & & & & & \\
				\hline
				12($\phi$7.5mm) & & & & & & & & & \\
				\hline
				13($\phi$15mm) & & & & & & & & & \\
				\hline
			\end{tabular}
		\end{table}
		
	\end{enumerate}
	
	\textbf{实验步骤：}
	
	\begin{enumerate}
		\item 实验前的准备：按照参考资料中的步骤，检查实验仪器，排除空气，调节回路水平，使实验管道内充满液体，且各测压管水位相同。
		\item 实验中的设置：将细软管与测压水柱相连，将止水夹夹住细软管，使测压水柱与大气隔开，防止水位波动。
		\item 实验前的状态检查：检查流量计是否清零，检查电测流装置是否正常工作，检查测压水柱是否与大气隔开。
		\item 通过调节泵转速，设置不同的流量条件：按照实验方案中的控制范围，选择三种不同的泵转速，使流量分别为高、中、低三种情况，例如，5 L/min、3 L/min、1 L/min。
		\item 实验记录：当流量稳定后，记录流量计的读数，观测各测压管的水位，记录数据。将数据填入实验记录表中。重复步骤4和步骤5，直到完成三种流量条件下的实验记录。
	\end{enumerate}
	
	
	\subsection{参考资料}
	\begin{enumerate}
		\item 伯努利原理
		
		\textbf{跨过流线的运动方程:}
		\[ \frac{\partial \rho}{\partial t} + \nabla \cdot (\rho \mathbf{v}) = 0 \]
		这个方程表示流体的质量守恒，即在任何封闭的控制体内，流入和流出的质量之差等于质量的变化率。
		
		\textbf{沿着流线的运动方程:}
		\[ \frac{D \mathbf{v}}{D t} = -\frac{1}{\rho} \nabla p + \mathbf{g} + \nu \nabla^2 \mathbf{v} \]
		这个方程表示流体的动量守恒，即在任何封闭的控制体内，外力和内力之和等于动量的变化率。其中，$\frac{D}{D t}$是物质导数，表示沿着流体微团的运动方向的变化率，$\nabla p$是压力梯度，表示压力的空间变化，$\mathbf{g}$是重力加速度，表示重力作用，$\nu$是运动粘度系数，表示流体的粘滞性，$\nabla^2 \mathbf{v}$是速度的拉普拉斯算子，表示速度的空间变化的变化率。
		
		\textbf{伯努利定律:}
		\[ p + \frac{1}{2} \rho v^2 + \rho g z = \text{constant} \]
		这个方程表示在理想流体（无粘性、无热传导、无压缩性）的稳定流动中，沿着流线的任意两点，流体的静压力、动压力和重力势能之和是一个常数。这个方程可以用来分析流体的速度、压力和高度之间的关系，例如水流从高处向低处流动时，速度会增加，压力会减小。
		
		\textbf{对于可压缩的流体:}
		\[ \frac{\partial}{\partial t} \left( \frac{1}{2} \rho v^2 + \rho e + \rho \Phi \right) + \nabla \cdot \left( \frac{1}{2} \rho v^2 \mathbf{v} + \rho e \mathbf{v} + p \mathbf{v} + \rho \Phi \mathbf{v} \right) = 0 \]
		这个方程表示在可压缩流体（密度随压力和温度变化）的流动中，流体的动能、内能和势能之和是一个守恒量。其中，$e$是比内能，表示单位质量的流体的内能，$\Phi$是保守力的势函数，表示单位质量的流体的势能，例如重力势能为$gz$，电势能为$q\phi$，其中$q$是电荷密度，$\phi$是电势。
		
		\textbf{在工程领域:}
		\[ \frac{\partial}{\partial t} \left( \frac{1}{2} \rho v^2 + \rho e \right) + \nabla \cdot \left( \frac{1}{2} \rho v^2 \mathbf{v} + \rho e \mathbf{v} + p \mathbf{v} \right) = 0 \]
		这个方程表示在可压缩流体的绝热流动中，流体的动能和内能之和是一个守恒量。这个方程忽略了保守力的势能项，因为在高海拔的地方，重力势能的变化相对于动能和内能的变化很小，而且流体流动的时间很短，不会受到其他保守力的影响。
		
		\textbf{在可压缩流体可以应用的地方:}
		\[ p + \frac{1}{2} \rho v^2 = \text{constant} \]
		这个方程表示在可压缩流体的等熵流动中，沿着流线的任意两点，流体的静压力和动压力之和是一个常数。这个方程忽略了重力势能项，因为在可压缩流体的流动中，高度变化很小，不会对流体的压力和速度产生显著的影响。
		
		\textbf{另一个适合使用在热力学的公式:}
		\[ h + \frac{1}{2} v^2 + gz = \text{constant} \]
		这个方程表示在可压缩流体的等焓流动中，沿着流线的任意两点，流体的比焓、动能和重力势能之和是一个常数。其中，$h$是比焓，表示单位质量的流体的焓，等于比内能加上比压力功，即$h = e + \frac{p}{\rho}$。这个方程可以用来分析流体的温度、压力和速度之间的关系，例如气体从高温高压的容器中喷出时，温度会降低，速度会增加。
		
		\item 层流
		
		\textbf{层流的雷诺数（Reynolds number）}是一个无量纲数，用来判断流体流动的状态是层流还是湍流。它的定义是
		\[ Re=\frac{\rho v L}{\mu} \]
		其中 $\rho$ 是流体密度，$v$ 是流体速度，$L$ 是特征长度，$\mu$ 是流体动力粘度。一般来说，当雷诺数小于某个临界值时，流体流动为层流，当雷诺数大于某个临界值时，流体流动为湍流。
		
		\textbf{层流的哈根-普瓦萨尔方程（Hagen–Poiseuille equation）}是一个描述层流管流的方程，它给出了管内流体的流量和压力降之间的关系。它的表达式是
		\[ Q=\frac{\pi r^4 \Delta p}{8 \mu L} \]
		其中 $Q$ 是流量，$r$ 是管道半径，$\Delta p$ 是管内两端的压力差，$\mu$ 是流体动力粘度，$L$ 是管道长度。
		
		\textbf{层流的布朗特-瓦西尔方程（Blasius boundary layer equation）}是一个描述层流边界层的方程，它给出了沿平板方向的速度分布。它的形式是
		\[ f''' + \frac{1}{2} f f'' = 0 \]
		其中 $f$ 是一个无量纲函数，与流体速度和平板坐标有关。这个方程是一个非线性的常微分方程，没有解析解，只能用数值方法求解。
		
		\item 湍流
		
		\textbf{湍流的雷诺数（Reynolds number）}是一个无量纲数，用来判断流体流动的状态是层流还是湍流。它的定义是
		\[ Re=\frac{\rho v L}{\mu} \]
		其中 $\rho$ 是流体密度，$v$ 是流体速度，$L$ 是特征长度，$\mu$ 是流体动力粘度。一般来说，当雷诺数大于某个临界值时，流体流动为湍流，当雷诺数小于某个临界值时，流体流动为层流。
		
		\textbf{湍流的科尔莫戈洛夫理论（Kolmogorov theory）}是一个描述湍流中能量在不同尺度上转移的理论，它提出了三个基本假设和三个基本尺度。它的核心公式是
		\[ \epsilon = \frac{u^3}{L} \]
		其中 $\epsilon$ 是湍流的能量耗散率，$u$ 是湍流的速度涨落，$L$ 是湍流的积分尺度。
		
		\textbf{湍流的纳维-斯托克斯方程（Navier–Stokes equations）}是一组描述流体运动的偏微分方程，它们是牛顿运动定律在流体上的应用。它们的一般形式是
		\[ \rho \left(\frac{\partial \mathbf{u}}{\partial t} + \mathbf{u} \cdot \nabla \mathbf{u}\right) = -\nabla p + \mu \nabla^2 \mathbf{u} + \mathbf{f} \]
		其中 $\rho$ 是流体密度，$\mathbf{u}$ 是流体速度，$p$ 是流体压力，$\mu$ 是流体动力粘度，$\mathbf{f}$ 是流体受到的外力。
		
		\item 皮托管
		\begin{itemize}
			\item \textbf{皮托管的原理和用途：} 皮托管是一种利用流体的压强差来测量流体的运动速度的仪器，由法国工程师亨利·皮托发明并由亨利·达西改进。皮托管常用于测量飞行器的空速和工业设施中的气体的流动速度。
			\item \textbf{皮托管的结构和测量方法：} 皮托管由一个直接处于流体中的管道组成，该管道中的流体停滞，测量其压强为滞压或总压。与机身侧面的静压孔测得的静压相减，得到动压，再通过空速计转换为表速。皮托管还可以与静压孔合并为皮托静压管，外管用来测量静压，内管用来测量总压。
			\item \textbf{皮托管的加热和失效：} 皮托管通常带有加热设备以防止皮托管被冰堵塞。皮托管失效会造成灾难性后果，如X-31、秘鲁航空603号班机、伯根航空301号班机和法国航空447号航班等事故。
			\item \textbf{皮托管的局限性：} 皮托管只能测量某一点的局部速度而不是整条管线的平均速度。皮托管也不适用于测量一些流体的速度，如马赫空速表。
		\end{itemize}		
		
	\end{enumerate}
%	\subsection{实验前思考题}

	
	
	
	%实验记录	
	\clearpage
	\begin{table}
		\renewcommand\arraystretch{1.7}
		\centering
		\begin{tabularx}{\textwidth}{|X|X|X|X|}
			\hline
			专业： & 物理学 & 年级： & 2022级 \\
			\hline
			姓名： & 杨舒云 & 学号： & 22344020\\
			\hline
			室温： & 23摄氏度 & 实验地点： & 实验室509 \\
			\hline
			学生签名：& 杨舒云 & 评分： &\\
			\hline
			实验时间：& 2023/11/30 & 教师签名：&\\
			\hline
		\end{tabularx}
	\end{table}
	
	\section{LabX4 伯努利原理实验 \quad\heiti 实验记录}
	\subsection{实验内容、步骤与结果}
	\subsubsection{实验前的准备（一般已准备好）}
	\begin{enumerate}
		\item 关闭出水阀门；
		\item 打开进水阀门，按下流量显示仪上的离心泵开关，启动离心泵，向实验回路供水；
		\item 在实验回路接近放满时，调节进水阀门，使实验回路达到溢流水平，并保持有一定的溢流；
		\item 适度打开出水阀门，使实验回路出流，当实验回路完全充满液体后，且12根测压管水位稳定且保持一致后，关闭出水阀门和进水阀门；
		\item 在作上述准备工作的过程中，应排尽管路和测压管中的空气，逐步排气和注水；
		\item 测量前，应仔细检查并调节电测流装置，使其能够正常工作。
		\item 调节回路水平（通过气泡判断），使回路位置（z）相等。		
	\end{enumerate}
	
	\clearpage
	\subsubsection{实验}
	\textbf{合作者：戴鹏辉22344016}。
	\begin{enumerate}
		\item 调节泵转速，当流量计读数稳定后，记录流量值（注意需要转换为国际单位），观测各测压管读数，记录数据。记录数据后，再次调节泵转速，分别记录高、中、低三种流量条件下的水头变化。这里的三种状态均应该保持测压管中有数据可读（即有压流动状态）；
		\item 选择一次大流量情况下，绘制测压管水头线和总水头线。将以上测量数据记入实验记录表（\cref{tab:tab1}）；
		
		\begin{table}[h]
			\centering
			\resizebox{\textwidth}{!}{%
			\begin{tabular}{|c|c|c|c|c|c|c|c|c|c|c|c|c|c|}
				\hline
				流量/($m^3/h$) & 1($\phi$20mm) & 2($\phi$20mm) & 3($\phi$15mm) & 4($\phi$15mm) & 5($\phi$15mm) & 6($\phi$15mm) & 7($\phi$15mm) & 8($\phi$15mm) & 9($\phi$15mm) & 10($\phi$15mm) & 11($\phi$15mm) & 12($\phi$7.5mm) & 13($\phi$15mm) \\ \hline
				0    & 560   & 559.8 & 560   & 559   & 560   & 559.8 & 559.8 & 560   & 560   & 559   & 558   & 558.6 & 560   \\ \hline
				0.113 & 568   & 568.5 & 565.5 & 566.2 & 562.9 & 562.5 & 564.8 & 561.5 & 561.5 & 563.5 & 556.9 & 527.7 & 541   \\ \hline
				0.15 & 573   & 573.5 & 570   & 572   & 565   & 564.5 & 569.5 & 562.5 & 561.5 & 566.5 & 556   & 504   & 526.5 \\ \hline
				0.2  & 583   & 586   & 577   & 582   & 570   & 568.8 & 575.8 & 565   & 563.8 & 571.5 & 556.8 & 460.5 & 508.5 \\ \hline
				0.25 & 595   & 599   & 585   & 594   & 574   & 573   & 584.5 & 567.5 & 565.5 & 578   & 557   & 403   & 489   \\ \hline
				0.3  & 609.8 & 615.5 & 595   & 609   & 580   & 578.5 & 594.5 & 570.5 & 568.5 & 586   & 555   & 336   & 465   \\ \hline
				0.35 & 626   & 634   & 606   & 626   & 586   & 584   & 606.5 & 574   & 570   & 595.2 & 552   & 265   & 431   \\ \hline
				0.4  & 645   & 656   & 619   & 645   & 593.8 & 592   & 620   & 577.2 & 572.7 & 605.5 & 549   & 185   & 390.5 \\ \hline
			\end{tabular}
		}
			\caption{实验记录表}
			\label{tab:tab1}
		\end{table}
		
		\item 作图展示（\cref{fig:fig1}）
		
		\begin{figure}[htbp]
			\centering
			\includegraphics[width=0.7\textwidth]{LabX4Gra1.png}
			\caption{测压管水头线和总水头线绘制图}
			\label{fig:fig1}
		\end{figure}
		
		数据分析见\textbf{实验数据分析}部分。
	\end{enumerate}
	

	
	
	\clearpage
	\subsection{原始数据记录}
	实验记录本上的原始数据见\cref{fig:figAD1}（两人分工情况为：记录数据： 戴鹏辉；调整实验仪器与读数： 杨舒云）。
	
	\begin{figure}[htbp]
		\centering
		\subfloat[]{
			\includegraphics[width=0.4\textwidth]{LabX4GraAD1.jpg}
		}
		\subfloat[]{
			\includegraphics[width=0.4\textwidth]{LabX4GraAD2.jpg}
		}
		\caption{原始数据}
		\label{fig:figAD1}			
	\end{figure}
	
	实验台桌面整理见\textbf{附件}部分(\cref{fig:figAD3})。
	
%	\subsection{实验过程中遇到的问题记录}
%	\begin{enumerate}
%		\item 
%	\end{enumerate}
	
	
	
	%分析与讨论	
	\clearpage
	\begin{table}
		\renewcommand\arraystretch{1.7}
		\begin{tabularx}{\textwidth}{|X|X|X|X|}
			\hline
			专业：& 物理学 &年级：& 2022级\\
			\hline
			姓名： & 杨舒云 & 学号：& 22344020\\
			\hline
			日期：& 2023/11/30 & 评分： &\\
			\hline
		\end{tabularx}
	\end{table}
	
	\section{LabX4 伯努利原理实验 \quad\heiti 分析与讨论}
	\subsection{实验数据分析（用数据回答问题）}
	
	\subsubsection{速度计算、分析}
	\begin{enumerate}
		\item 能量方程实验管上4组测压管的任一组都相当于一个毕托管，可测得管内轴心的流体速度：$v=\sqrt{2g\Delta h}$，其中 Δh为相应截面上两侧压管的水头差（即动压水头）；计算管内的平均流速$u=\dfrac{4Q}{πD^2}$，与v的数据比较。
		
		分别计算得到如\cref{tab:tab2}所示结果。
		
		\begin{table}[h]
			\centering
			\begin{tabular}{|c|c|c|c|c|c|c|c|c|c|}
				\hline
				流量/($m^3/h$) & & 0.000 & 0.113 & 0.150 & 0.200 & 0.250 & 0.300 & 0.350 & 0.400 \\
				\hline
				轴心流体速度v & A & 0.000 & 0.099 & 0.099 & 0.242 & 0.280 & 0.334 & 0.396 & 0.464 \\
				& B & 0.000 & 0.117 & 0.198 & 0.313 & 0.420 & 0.524 & 0.626 & 0.714 \\
				& C & 0.000 & 0.203 & 0.305 & 0.354 & 0.464 & 0.547 & 0.649 & 0.729 \\
				& D & 0.000 & 0.198 & 0.297 & 0.373 & 0.475 & 0.569 & 0.674 & 0.774 \\
				\hline
				平均流速u      & A & 0.000 & 0.100 & 0.133 & 0.177 & 0.221 & 0.265 & 0.309 & 0.354 \\
				& B & 0.000 & 0.178 & 0.236 & 0.314 & 0.393 & 0.472 & 0.550 & 0.629 \\
				& C & 0.000 & 0.178 & 0.236 & 0.314 & 0.393 & 0.472 & 0.550 & 0.629 \\
				& D & 0.000 & 0.178 & 0.236 & 0.314 & 0.393 & 0.472 & 0.550 & 0.629 \\
				\hline
			\end{tabular}
			\caption{流速比较表}
			\label{tab:tab2}
		\end{table}					
		
		\item 为什么管轴心流速总是大于管内平均流速u？层流和湍流条件下，管中心流速和管内平均流速的差异如何变化？请给出分析，猜测管内的流速是如何分布的。
		
		将实验数据导入origin，绘制不同测压点的轴心流体速度与平均流速对比图（\cref{fig:fig2}、\cref{fig:fig4}）。
		
		\begin{figure}[htbp]
			\centering
			\subfloat[]{
				\includegraphics[width=0.5\textwidth]{LabX4Gra2.jpg}
			}
			\subfloat[]{
				\includegraphics[width=0.5\textwidth]{LabX4Gra3.jpg}
			}
			\caption{轴心流体速度与平均流速对比图}
			\label{fig:fig2}			
		\end{figure}
		
		\begin{figure}[htbp]
			\centering
			\subfloat[]{
				\includegraphics[width=0.5\textwidth]{LabX4Gra4.jpg}
			}
			\subfloat[]{
				\includegraphics[width=0.5\textwidth]{LabX4Gra5.jpg}
			}
			\caption{轴心流体速度与平均流速对比图}
			\label{fig:fig4}			
		\end{figure}
		
		通过对比图，我们可以认为轴心流速通常大于管内平均流速。下面做一些具体解释。
		
		管轴心流速总是大于管内平均流速的原因可以从流体力学的角度进行分析。在一个圆管中，流体的流速分布不是均匀的，而是呈现出剖面速度分布，这是因为管壁对流体有阻力作用。在层流条件下，这种阻力导致流体在管壁附近的流速远小于管中心的流速。流体速度沿着从管壁到管中心的径向逐渐增加，达到最大值，即管轴心流速。
		
		在层流条件下，流体的流动是有序的，各层流体平滑且互不扰动地滑过相邻层，速度分布呈抛物线状，这意味着管中心的流速最大，而管壁处的流速为零。因此，管中心的流速（轴心流体速度）远大于平均流速。平均流速是横截面上所有不同流速的平均值。
		
		湍流条件下，流体的流动是无序的，各层流体间发生混合和涡流，速度分布更加均匀，但管中心的流速仍然是最大的。与层流相比，湍流的速度分布较为平坦，但中心仍然高于边缘。虽然湍流导致更多的动能传递到管壁附近，从而提高了那里的流速，但管中心的流速仍然是最大的。因此，即便在湍流条件下，管轴心流速也大于平均流速，只是两者之间的差异不如层流条件下那么显著。
		
		管内的流速分布如下：在层流条件下，速度分布呈现出明显的抛物线状，中心高、边缘低。在湍流条件下，速度分布更加均匀，但仍然可以观察到从管壁到管中心的递增趋势。总的来说，我们可认为在管道的中心，速度较大，而靠近管道壁的地方，速度较小。
	\end{enumerate}
	
	\subsubsection{能量损耗分析：总水头变化，摩擦压降}
	\begin{enumerate}
		\item 在测量所得实验数据基础上，计算出伯努利方程试验管各取压点截面的相应动压水头和压强水头（\cref{tab:tab3}）。
		
		\begin{table}[h]
			\centering
			\begin{tabular}{|c|c|c|c|c|c|c|c|c|}
				\hline
				流量/($m^3/h$) & 0 & 0.113 & 0.15 & 0.2 & 0.25 & 0.3 & 0.35 & 0.4 \\
				\hline
				动压水头A & -0.2 & 0.5 & 0.5 & 3 & 4 & 5.7 & 8 & 11 \\
				B & -1 & 0.7 & 2 & 5 & 9 & 14 & 20 & 26 \\
				C & -0.1 & 2.1 & 4.75 & 6.4 & 11 & 15.25 & 21.5 & 27.1 \\
				D & -1 & 2 & 4.5 & 7.1 & 11.5 & 16.5 & 23.2 & 30.55 \\
				\hline
				压强水头A & 560 & 568 & 573 & 583 & 595 & 609.8 & 626 & 645 \\
				B & 560 & 565.5 & 570 & 577 & 585 & 595 & 606 & 619 \\
				C & 559.9 & 562.7 & 564.75 & 569.4 & 573.5 & 579.25 & 585 & 592.9 \\
				D & 560 & 561.5 & 562 & 564.4 & 566.5 & 569.5 & 572 & 574.95 \\
				\hline
			\end{tabular}
			\caption{动压水头和压强水头表}
			\label{tab:tab3}
		\end{table}
				
		\item 绘制一定工况下的4个取压点截面的各种水头和总水头的水头线。
		
		\cref{fig:fig6}分别给出了动压水头（左）、压强水头（中）和总水头（右）的绘制曲线。
		
		\begin{figure}[htbp]
			\centering
			\subfloat[]{
				\includegraphics[width=0.33\textwidth]{LabX4Gra6.jpg}
			}
			\subfloat[]{
				\includegraphics[width=0.33\textwidth]{LabX4Gra7.jpg}
			}
			\subfloat[]{
				\includegraphics[width=0.33\textwidth]{LabX4Gra8.jpg}
			}
			\caption{各种水头线}
			\label{fig:fig6}			
		\end{figure}
		
		\item 运用伯努利方程进行分析，解释各水头的变化规律，例如：
		\begin{enumerate}
			\item 可以看出能量损失沿着流体流动的方向是如何变化的？测压管水头线和总水头线的沿程变化趋势有何不同？为什么？
			
			根据总水头曲线，我们可以看出总水头随着流体流动方向逐渐降低，即能量随着流体流动方向逐渐降低。
			
			测压管水头线是沿水流方向各个测点的测压管液面的连线，它反映的是流体的势能。测压管水头线沿程可能下降，也可能上升（当管径沿流向增大时），因为管径增大时流速减小，动能减小而压能增大，如果压能的增大大于水头损失时，水流的势能就增大，测压管水头就上升。总水头线是在测压管水头线的基线上再加上流速水头，它反映的是流体的总能量。
			由于沿流向总是有水头损失，所以总水头线沿程只能下降，不能上升。
			
			\item 随着流量增加，测压管水头线变化趋势如何变化？
			为什么？请计算摩擦压降（压力损失），它与管径、流速和流动距离有怎样的关系？（达西公式）。
			
			对于动压水头的变化，由流体连续性方程可知，管径减小，流速会增加。对应于B、C、D处的管径都是15mm，而A点处的管径是20mm。因此绘制曲线表面B、C、D点的流速差距不大，而A点与三者差距较大。 
			
			随着流量增加，测压管水头线的变化趋势是下降的。这是因为流量和测压管水头线的坡度成正比，而测压管水头线的坡度等于管道起端的测压管水头减去末端的测压管水头。通常管道的起端的测压管水头是基本不变的，所以末端的测压管水头越低，测压管水头线坡度就越大，流量也就越大。
			
			下面讨论摩擦压降相关：
			
			计算摩擦压降如\cref{tab:tab4}所示：
			\begin{table}[h]
				\centering
				\begin{tabular}{|c|c|c|c|c|c|c|c|c|}
					\hline
					流量/($m^3/h$) & 0 & 0.113 & 0.15 & 0.2 & 0.25 & 0.3 & 0.35 & 0.4 \\
					\hline
					摩擦压降AB & 0.8 & 2.3 & 1.5 & 4 & 5 & 6.5 & 8 & 11 \\
					BC & -0.8 & 1.4 & 2.5 & 6.2 & 9.5 & 14.5 & 19.5 & 25 \\
					CD & 0.8 & 1.3 & 3 & 4.3 & 6.5 & 8.5 & 11.3 & 14.5 \\
					\hline
				\end{tabular}
				\caption{摩擦压降计算表}
				\label{tab:tab4}
			\end{table}
			
			按实验原理部分给出的达西公式$\Delta p = \lambda \dfrac{l}{d} \dfrac{1}{2} \rho v^2$，
			流体密度越大，速度越快，摩擦压降越大；
			管道内径越小，长度越长，摩擦压降越大；
			管道的表面粗糙度越大，摩阻系数越大，摩擦压降越大。
			
			绘制总水头线的摩擦压降与流量的关系图来分析各个因素对摩擦压降的影响，如\cref{fig:fig9}所示。
			
			\begin{figure}[htbp]
				\centering
				\includegraphics[width=0.7\textwidth]{LabX4Gra9.jpg}
				\caption{摩擦压降-流量关系曲线}
				\label{fig:fig9}
			\end{figure}
			
			对于流速来说，按达西公式，流体的摩擦压降正比于流速的平方，而流量与流速成正比，则流体的摩擦压降与流量的平方成正比。选取BC段的数据，利用Origin进行拟合，得到如\cref{fig:fig10}所示结果。可以发现，拟合效果较好，即达西公式对BC段比较成立。
			
			\begin{figure}[htbp]
				\centering
				\includegraphics[width=0.7\textwidth]{LabX4Gra10.jpg}
				\caption{摩擦压降-流量关系拟合曲线}
				\label{fig:fig10}
			\end{figure}
			
			对于管径来说，由达西公式可知，摩擦压降与管径成反比。AB之间的距离与CD之间的距离相同，AB段管径一半为20mm，一半为15mm，而CD段全部为15mm。因此理论上CD段的摩擦压降更大，通过图像可以看出，这与实验得到的结果是相吻合的，这至少在定性上说明了达西公式的成立。
			
			\item 实验结果是否符合连续方程（质量守恒）？即当流量一定时，管径与当地流速有怎样的关系?
			
			由连续性方程计算并作图得到\cref{tab:tab5}和\cref{fig:fig11}所示结果。
			
			\begin{table}[h]
				\centering
				\begin{tabular}{|c|c|c|c|c|}
					\hline
					流量 & 计算流量值A & B & C & D \\
					\hline
					0 & 0 & 0 & 0 & 0 \\
					0.113 & 3.11002E-05 & 2.0699E-05 & 3.58517E-05 & 3.49877E-05 \\
					0.15 & 3.11002E-05 & 3.49877E-05 & 5.39197E-05 & 5.24816E-05 \\
					0.2 & 7.61796E-05 & 5.53204E-05 & 6.25879E-05 & 6.59219E-05 \\
					0.25 & 8.79646E-05 & 7.42201E-05 & 8.20534E-05 & 8.38976E-05 \\
					0.3 & 0.000105006 & 9.25688E-05 & 9.6613E-05 & 0.000100495 \\
					0.35 & 0.000124401 & 0.000110641 & 0.000114715 & 0.000119164 \\
					0.4 & 0.000145873 & 0.00012615 & 0.000128791 & 0.000136743 \\
					\hline
				\end{tabular}
				\caption{计算流量值}
				\label{tab:tab5}
			\end{table}
			
			\begin{figure}[htbp]
				\centering
				\includegraphics[width=0.7\textwidth]{LabX4Gra11.jpg}
				\caption{计算流量值-流量关系曲线}
				\label{fig:fig11}
			\end{figure}
			
			根据上述结果，我么可以认为A、B、C和D四管的流量计算值几乎相同，即实验结果符合连续方程。
			
			除此之外，我么可以得到$V = \frac{4Q}{\pi d^2}$，即当流量一定时，管径与当地流速成反比（管径越大，当地流速越小。管径越小，当地流速越大。）			
		\end{enumerate}		
	\end{enumerate}
	
	\subsubsection{伯努利原理：文丘里管}
	\begin{enumerate}
		\item 按文丘里管管径计算出E、F、G点对应截面流速；
		
		按文丘里管中流速计算公式$v=\sqrt{\dfrac{2(p_1-p_2)}{\rho(\frac{S_1^2}{S_2^2}-1)}}$，计算得到\cref{tab:tab6}所示结果。
		
		\begin{table}[h]
			\centering
			\begin{tabular}{|c|c|c|c|c|c|c|c|c|}
				\hline
				流量 $m^3/h$ & 0.000 & 0.113 & 0.150 & 0.200 & 0.250 & 0.300 & 0.350 & 0.400 \\
				\hline
				截面流速 E & 0.000 & 0.062 & 0.083 & 0.113 & 0.143 & 0.171 & 0.196 & 0.220 \\
				F & 0.055 & 0.168 & 0.219 & 0.320 & 0.428 & 0.525 & 0.595 & 0.662 \\
				G & 0.014 & 0.042 & 0.055 & 0.080 & 0.107 & 0.131 & 0.149 & 0.166 \\
				\hline
			\end{tabular}
			\caption{截面流速}
			\label{tab:tab6}
		\end{table}
		
		\cref{fig:fig12}展示了截面流速-流量的可视化关系。
		
		\begin{figure}[htbp]
			\centering
			\includegraphics[width=0.7\textwidth]{LabX4Gra12.jpg}
			\caption{截面流速-流量关系曲线}
			\label{fig:fig12}
		\end{figure}
		
		\item 计算不同管径取压点的水柱高度差，它与当地流速有怎样的关系？
		
		通过计算后绘制相邻两测压点的水柱高度差-流量关系曲线，如\cref{fig:fig13}所示。
		
		\begin{figure}[htbp]
			\centering
			\includegraphics[width=0.7\textwidth]{LabX4Gra13.jpg}
			\caption{水柱高度差-流量关系曲线}
			\label{fig:fig13}
		\end{figure}
		
		根据伯努利定理，流体在管道中的压力和速度满足以下方程：$p + \frac{1}{2}\rho V^2 = C$。其中，$p$ 是流体的压力，$\rho$ 是流体的密度，$V$ 是流体的速度，$C$ 是一个常数。在文丘里管的入口截面和最小截面处，分别取压点，测量水柱高度差，即：$\Delta h = h_1 - h_2$。其中，$h_1$ 是入口截面处的水柱高度，$h_2$ 是最小截面处的水柱高度。假设流体的密度为 $\rho$，水的密度为 $\rho_w$，重力加速度为 $g$，则水柱高度差与流体压力差的关系为：$\Delta p = \rho_w g \Delta h$。代入伯努利方程，得到：$\frac{1}{2}\rho V_1^2 - \frac{1}{2}\rho V_2^2 = \rho_w g \Delta h$。其中，$V_1$ 是入口截面处的流体速度，$V_2$ 是最小截面处的流体速度。如果假设流体是不可压缩的，即 $\rho$ 是一个常数，那么流体的流量 $Q$ 也是一个常数，即：$Q = A_1 V_1 = A_2 V_2$。其中，$A_1$ 是入口截面的面积，$A_2$ 是最小截面的面积。由此，可以解出：$V_1 = \frac{A_2}{A_1} V_2$。代入水柱高度差的方程，得到：$\Delta h = \frac{\rho}{2\rho_w g} \left( \frac{A_1^2 - A_2^2}{A_2^2} \right) V_2^2$。从这个方程可以看出，水柱高度差与最小截面处的流体速度的平方成正比（当地流速越大，水柱高度差越大；当地流速越小，水柱高度差越小）。同时，水柱高度差也与管道的入口截面和最小截面的面积有关入口截面越大（最小截面越小，水柱高度差越大；入口截面越小，最小截面越大，水柱高度差越小）。
		
		对比我们实验得到的结果：相邻水柱高度差与流量成平方关系，即相邻水柱高度差与当地流速成平方关系。		
		
		\item 数据是否可以用伯努利方程的描述：如何处理11-12点和12-13点的差异？
		能否通过11点与13点的水头差异，扣除能量损失的影响？	
		
		伯努利方程是描述流体流动中能量守恒的一个重要方程，它考虑了流体在运动过程中的位能、压能和动能的转换。在实验中，可以通过测量不同截面的压强和速度来应用伯努利方程。伯努利方程中涉及的能量损失（如摩擦损失）通常在实际流体中需要额外考虑。伯努利方程在考虑能量损失的情况下可以表示为：\[ z_1 + \frac{p_1}{\rho g} + \frac{v_1^2}{2g} = z_2 + \frac{p_2}{\rho g} + \frac{v_2^2}{2g} + h_{1-2} \] 其中 \( h_{1-2} \) 为能量损失项。因此，通过比较11点和13点的水头差异，确实可以估计在这两点之间的能量损失。
		
		具体来看，EF和FG 的水柱高度差之间存在差异，这是因为随着水流方向，总能量在逐渐降低。即EF与FG的水柱高度差的差值就是能损失的能量。	
	\end{enumerate}
	
	\subsection{用实验数据回答的思考题}
	\begin{enumerate}
		\item 大气压强是否会对实验结果有影响？
		
		在伯努利方程中，大气压强会影响流体的压强水头，当大气压强增加时，压强水头也会增加。因此在伯努利原理实验中，大气压强可能对实验结果产生影响，尤其是在涉及到开放式测压管的情况下。大气压强作用于液体表面，会影响测量的压强值。如果实验涉及到流体中的速度和压力变化，大气压的变化可能导致实验结果的偏差。
		
		但是，实际在伯努利方程中，通常考虑的是相对压强，即测量点的压强与大气压强的差值。因此，只要实验过程中大气压强保持相对稳定，其影响可以被相对化处理。在本实验中， 并未考虑大气压强的变化量，可认为其是一个常数，其对实验结果的影响很小。
		
		\item 如未充分排气，请分析实验过程中排气对于实验结果的影响。
		\begin{enumerate}
			\item 影响流速分布: 伯努利方程假定了流体是连续和均匀的，未充分排气可能导致管道或容器中存在气泡或空腔，这些空腔可能干扰流体的流动，导致流速分布不均匀。
			\item 增加摩擦: 空气困留在管道或容器中可能会增加摩擦阻力，从而影响流体的动能 和压力分布。
			\item 误差和不确定性: 未充分排气可能引入不确定性和误差，使实验结果不可靠。
		\end{enumerate}
		总得来说，若实验管路中未充分排除空气，可能会导致实验数据的不准确。空气泡会影响流量计的读数，导致流速计算不准确。此外，空气泡在流动过程中可能会被压缩或膨胀，这会影响液体的压强和速度测量，从而影响到伯努利方程的应用。
		
				
	\end{enumerate}
	

%	\subsection{实验后思考题}
	
	
	
	
	\clearpage
	\section{LabX4 伯努利原理实验 \quad\heiti 参考文献与附件}
	\subsection{参考文献}
	[1] 维基百科
	
	[2] 实验讲义
	
	\subsection{附件}
	\begin{figure}[htbp]
		\centering
		\includegraphics[width=0.7\textwidth]{LabX4GraAD3.jpg}
		\caption{实验台桌面整理}
		\label{fig:figAD3}
	\end{figure}
	
	
	
\end{document}